\documentclass{article}
\usepackage[left=2cm, right=2cm, top=2cm, bottom=2cm]{geometry}
\usepackage{graphicx}
\usepackage{amsmath}
\usepackage{amssymb}
\usepackage{amsthm}
\usepackage{fancyhdr}
\usepackage{verbatim}
\usepackage{listings}
\usepackage{xcolor}
\usepackage{pdfpages}
\usepackage{cancel}
\usepackage{physics}
\usepackage{esint}

\lstset{
    language=C++,
    basicstyle=\ttfamily\footnotesize,
    keywordstyle=\color{blue},
    commentstyle=\color{green},
    stringstyle=\color{red},
    numbers=left,
    numberstyle=\tiny\color{gray},
    frame=single,
    breaklines=true,
    captionpos=b,
}


\title{MTH 553: Lecture \# 38}
\author{Cliff Sun}

\newtheorem{theorem}{Theorem}[section]
\newtheorem{lemma}[theorem]{Lemma}
\newtheorem{definition}[theorem]{Definition}
\newtheorem{conjecture}[theorem]{Conjecture}
\newtheorem{proposition}[theorem]{Proposition}
\newtheorem{corollary}[theorem]{Corollary}
\newtheorem{one minute paper}[theorem]{One Minute Paper}

\pagestyle{fancy}
\lhead{\textbf{\thepage}\ \ \nouppercase{\rightmark}}
\chead{MTH 553: Lecture \# 38}
\rhead{Cliff Sun}

\begin{document}

\maketitle

\section*{Lecture Span}
\begin{enumerate}
    \item Reaction diffusion traveling wave
\end{enumerate}

\section*{Traveling wave for a bistable reaction diffusion equation}

\begin{equation}
    u_t = u_{xx} + f(u)
\end{equation}

Here, $f(u)$ is a reaction term. This models the diffusion of a substance with chemical reaction generatin new material at the rate of $f(u)$. 

Suppose our curve $f$ has two humps whos middle intersects $a$. Note that the top hump is larger than the bottom hump. Here, $0 < a < 1$ and 

\begin{equation}
    \int_{0}^{1}f(z)dz > 0
\end{equation}

That is, there is more area over the axis than under the axis. But now, what does bistable mean? Here, $u \equiv 0$ is a stable solution of the PDE. 

If $u \equiv -\epsilon$ then $f(u) > 0$, so then $u$ approaches $0$. If $u \equiv \epsilon$, then $f(u) < 0$. So then $u$ also approaches $0$. This is to say that the solution $u \equiv 0$ is \underline{stable}.

Similarly, $u\equiv 1$ is also stable. However, $u = a$ is unstable. Therefore, we call the PDE bistable because there are two stable steady state solutions. 

Our goal is find 

\begin{equation}
    u(x,t) = v(x + ct)
\end{equation}

We need that 

\begin{equation}
    \begin{cases}
        v(-\infty) = 0 \\
        v'(\pm\infty) = 0
    \end{cases}
\end{equation}

Now, substitute into the PDE to obtain 

\begin{equation}
    cv'(s) = v''(s) + f(v(s))
\end{equation}

Where $s = x + ct \in \mathbb{R}$. Here, we will use the phase plane method to turn this 2nd order 
ODE into two first order ODEs. So then, we obtain 

\begin{align}
    v' &= w \\
    w' &= f(v) + cw
\end{align}

Then we want 

\begin{enumerate}
    \item $(v,w) \rightarrow (1,0)$ as $s \rightarrow \infty$
    \item $(v,w) \rightarrow (0,0)$ as $s \rightarrow -\infty$
\end{enumerate}

Here, to understand the system at $(0,0)$, we aim to linearize the ODE: 

\begin{equation}
    \begin{pmatrix}
        v'\\
        w'
    \end{pmatrix} = \begin{pmatrix}
        0 & 1 \\
        -f'(0) & c
    \end{pmatrix}\begin{pmatrix}
        v \\
        w
    \end{pmatrix}
\end{equation}

Then find the eigenvalues of $A$ (the above matrix) to obtain

\begin{equation}
    \lambda_{\pm} = \frac{c \pm \sqrt{c^2 - 4f'(0)}}{2}
\end{equation}

At $(1,0)$, the linearization is 

\begin{equation}
    \begin{pmatrix}
        (v - 1)'\\
        w'
    \end{pmatrix} = \begin{pmatrix}
         0 & 1 \\
         -f'(1) & c
    \end{pmatrix}\begin{pmatrix}
        v - 1\\
        w
    \end{pmatrix}
\end{equation}

Then the eigenvalues of 

\begin{equation}
    \lambda_{\pm} = \frac{c \pm \sqrt{c^2 - 4f'(1)}}{2}
\end{equation}

Note, $f'(0) < 0$ and $f'(1) < 0$, which then means that 

\begin{equation}
    \lambda_{-} < 0 < \lambda_{+}
\end{equation}

Therefore, you have an incoming eigenvector and an outgoing eigenvector. We claim that there is a unique $c >0$ such that the outgoing trajectories between $(0,0)$ and $(1,0)$ are connected. 

Here, to find such a $c$, use Lypunov function 

\begin{equation}
    E(v,w) = \int_{0}^vf(z)dz + \frac{1}{2}w^2
\end{equation}

Which increases along trajectories of ODE systems. Then 

\begin{align}
    \frac{d}{ds}E(v(s), w(s)) &= (f(v(s)) + v''(s))v'(s) \\
    &= cv'(s)^2 > 0
\end{align}

\end{document}